\documentclass[conference]{IEEEtran}
\usepackage[T1]{fontenc}
\usepackage{booktabs,multirow,amsmath,graphicx,xcolor,url}
% Double-anonymous review: \anontrue hides the author, the affiliation and the artifact location.
\newif\ifanon
\anonfalse

\begin{document}
\title{Beyond Grep: Intent-Aware Interception and At-the-Moment Teaching for Coding Agents\\(An Exploratory Study)}

\ifanon
\author{\IEEEauthorblockN{Anonymous author(s)}\IEEEauthorblockA{Anonymous institution}}
\else
\author{\IEEEauthorblockN{Hung T. Le}\IEEEauthorblockA{Independent researcher, Austin, Texas, USA \quad letronghung@gmail.com}}
\fi

\maketitle

\begin{abstract}
Coding agents explore a repository with text search, and code-intelligence tools offered beside that search go largely
unused: in 72 of our sessions agents never called a symbol lookup on their own. We study a different delivery: a hook
that intercepts the agent's own \texttt{grep}, classifies what the search asks (13 intent classes), asks a code graph
the question that answers it, and returns the answer beside or instead of the search, shaped like grep's output. We
report an exploratory study, not pre-registered. On 973 recorded agent searches, answers to a different question than
the search asked fell from 52 of 120 to none of 43 judged, and answer text from 150k to 41k characters. Over 24
implementation tickets (216 runs, one model) the hook raised the share of runs passing a hidden acceptance check and the
suite from 71\% to 79\%, and lowered cost per successful run by 7\%, at unchanged cost per run. One line appended to each
answer, naming the graph command that gives it, raised direct graph queries from 0.17 to 1.50 per run and lowered cost
per successful run by 19\% against the hook alone. Cost tracked turns ($r=0.92$), and two-thirds of all context was a
fixed load re-read every turn, so an answer is worth the turns it saves, not its size. We list the measurement errors we
made and corrected, and pre-specify the confirmatory study these results call for.
\end{abstract}

\begin{IEEEkeywords}
coding agents, LLM agents, code search, code graph, tool adoption, agent-computer interface, cost, empirical study
\end{IEEEkeywords}

\section{Introduction}
LLM coding agents~\cite{jimenez2024swebench,yang2024sweagent} navigate code with \texttt{grep}, file reads and line
ranges, and pay again for every read on each later turn. Code graphs~\cite{liu2024codexgraph,ouyang2025repograph} and
precomputed context~\cite{le2026precomputed} can answer the questions behind those searches more directly. But
availability is not adoption: an agent offered graph tools often does not call them. CodeCompass reports that 58\% of
trials with graph access made no graph call~\cite{paipuru2026codecompass}; in the 72 sessions behind this study, agents
called a symbol lookup zero times. Our earlier pre-registered study found that a graph pays when its content is
delivered in the prompt rather than behind tools~\cite{le2026precomputed}. Here we ask what to do in everyday use, where
nothing is precomputed for the task: \emph{deliver the graph's answer inside the tool call the agent already makes}.

Intercepting a search is safe only if the graph answers the \emph{same question}. A naive interceptor sees an identifier
and returns what the graph knows about it; we found that the earlier version of our hook did exactly that, and answered
a different question than the search asked 52 times in 120. This paper studies three mechanisms and their effect on
quality and cost:
\begin{itemize}
\item \textbf{Intent-preserving interception}: classify each search into one of 13 intents, map each intent to one graph
  question and an answer shaped like grep's output, and pass the search through untouched whenever equivalence is not
  established (Section~\ref{sec:system}).
\item \textbf{Teaching at the moment of need}: end each answer with the graph command that produced it, so that the next
  such question can go to the graph directly (Section~\ref{sec:teach}).
\item \textbf{A cost model for agent exploration}: turns, not answer size, drive cost, because a fixed load dominates the
  context re-read on every turn (Section~\ref{sec:cost}).
\end{itemize}
The evidence is exploratory: one model, one agent, tickets from the tool's own repository, three runs per arm. We report
it with its noise, the errors we made in measuring it (Section~\ref{sec:measure}), and the confirmatory design it
motivates (Section~\ref{sec:plan}).

\section{Related work}\label{sec:related}
\textbf{Agent-computer interfaces.} SWE-agent showed that the design of the agent's tools changes its success on
repository tasks~\cite{yang2024sweagent}. SWE-Edit separates inspection from editing into subagents with clean contexts
and reports both higher resolve rates and lower inference cost~\cite{zhang2026sweedit}; like us it attacks the context an
agent accumulates, by a different route. \textbf{Structure-aware search.} AutoCodeRover and SpecRover search code by
program structure and infer the intended behaviour of what they visit~\cite{zhang2024autocoderover,ruan2025specrover};
LocAgent and RepoGraph build graphs for localisation~\cite{chen2025locagent,ouyang2025repograph}. These systems
\emph{replace} the agent's search with their own; we keep the agent's own commands and add to them. \textbf{Graph
navigation and adoption.} CodeCompass serves structural dependencies over MCP and finds both a benefit on dependencies
with little lexical overlap and an adoption problem~\cite{paipuru2026codecompass}; it argues for explicit prompting. Our
teach line is a different intervention: no instruction up front, a one-line hint inside the answer the agent is already
reading. \textbf{Our earlier study} measured process-graph packs under pre-registration~\cite{le2026precomputed}; this
paper concerns the interception path, which needs no per-task precomputation.

\section{Intent-preserving interception}\label{sec:system}
The hook runs before each \texttt{Grep}, \texttt{grep}/\texttt{rg}/\texttt{git grep} command and \texttt{Read} call
in Claude Code, on a repository with a built graph (symbols, calls, references, string-key uses, tests). For each
search it:
\begin{enumerate}
\item \textbf{parses} the command part that searches (pipes, \texttt{-e} lists, \texttt{-A/-B/-C}, \texttt{-c},
  \texttt{-l}, \texttt{-v}, \texttt{--include}, globs, scope operands), and the grep dialect, so that \texttt{a|b} is an
  alternation only where grep reads it as one;
\item \textbf{classifies} the intent: read a body, locate a definition, list a file's symbols, call sites (plain,
  filtered by a pipe, in one file), count, list files, invert, literal text, non-code, unknown, not a search;
\item \textbf{asks} the graph the intent's question: a definition and its body, a file outline, or name uses with calls
  first; nothing for counts, file lists, inverted and literal searches;
\item \textbf{shapes} the answer as \texttt{file:line: text} lines with one wrapper line, within 1{,}500 characters, and
  adds nothing when its lines would only repeat grep's;
\item \textbf{decides}: \emph{replace} the search when the answer is complete and the graph is current;
  \emph{alongside} when it is useful but not complete; \emph{pass through} otherwise.
\end{enumerate}
The graph is treated as stale, and never replaces a search, when it was built at another commit or the searched scope
holds changed code files. The specification states falsifiable requirements: at most 5\% of answers to a different
question on a replay corpus, and no answer beside grep larger than grep's own output.

\section{Method}\label{sec:method}
\textbf{Replay.} 973 searches recorded in earlier benchmark sessions (Claude Code, Sonnet) are replayed through each hook
version on the same checkout. Answers were labelled \emph{same}, \emph{partly} or \emph{different} question against
the search, by rules first and an LLM judge after; disagreements were resolved by hand.

\textbf{Tri-mode benchmark.} For each of 24 implementation tickets of the 0.2.13 release (graph build, call
resolution, the hook itself, a triage daemon, the benchmark harness), three arms ran on the same commit: no graph
(\emph{none}), the previous release's hook (\emph{v0.2.12}), and the hook as it stood on the main branch when that
batch ran (\emph{main}), three runs each, Sonnet, a fresh checkout per run. A run passes when a hidden acceptance check,
removed from the agent's checkout and written before the batch, and the full test suite both pass. Each devkb arm took
its tool from an environment built at an exact commit, not from the checkout the agent edited.

\textbf{Direct-path arms.} On the last 6 tickets two more arms ran: the hook with the teach line (D1) and the hook with a
file outline on first open (D2), 18 runs per arm.

\textbf{Counting.} Turns are distinct assistant messages; tokens are the session totals the client reports; a graph
command counts as navigation when run before the agent's first edit of the tool's own code, and as a self-test after it
(the agents edit the repository that contains the tool).

\section{Results}\label{sec:results}
\subsection{Does the hook answer the question the search asks?}
\begin{table}[t]\centering\small
\caption{Replay of 973 recorded searches through each hook version.}\label{tab:replay}
\resizebox{\columnwidth}{!}{\begin{tabular}{lrrr}\toprule
Hook & Answered & Same question & Answer text\\\midrule
v0.2.12 (identifier $\rightarrow$ callers) & 120 & 34 (52 different) & 150k chars\\
+ one question per intent & 110 & 100 of 101 judged & 92k\\
+ grep-shaped, no repeats & 52 & 42 of 43 judged$^\ast$ & 41k\\\bottomrule
\multicolumn{4}{l}{\footnotesize $^\ast$ the 43rd was a judge artefact; 70 answers dropped as repeats of grep's lines.}
\end{tabular}}\end{table}
Table~\ref{tab:replay}: mapping each intent to its own question removed the wrong-question answers; shaping removed
answers that only repeated grep's output. The hook answers fewer searches, which is the intended outcome: the dropped
answers were noise.

\subsection{Quality and cost on implementation tickets}
\begin{table}[t]\centering\small
\caption{24 tickets, 72 runs per arm (Sonnet). Pass = hidden check and suite.}\label{tab:tri}
\resizebox{\columnwidth}{!}{\begin{tabular}{lrrrr}\toprule
Arm & Passed & Cost/run & Cost/pass & Turns vs none\\\midrule
none & 51 (71\%) & \$0.370 & \$0.522 & --\\
v0.2.12 hook & 58 (81\%) & \$0.383 & \$0.475 & $-5.2\%$\\
main hook & 57 (79\%) & \$0.383 & \$0.484 & $-4.7\%$\\\bottomrule
\end{tabular}}\end{table}
Both hook arms pass 6--7 more runs than no graph (Table~\ref{tab:tri}). Cost per run does not move: $+3.5\%$ on the
arithmetic mean, $-0.6\%$ on the geometric mean of per-ticket ratios; cost per successful run falls 7--9\%. With three
runs per arm a single ticket's difference of up to about 20\% is within noise; we read only the release-wide figures.
The current hook does not beat v0.2.12 on these tickets: the intent work made answers right and small, which shows on
the replay, not as more passing runs on small tickets.

\subsection{Teaching at the moment of need}\label{sec:teach}
\begin{table}[t]\centering\small
\caption{Last 6 tickets, 18 runs per arm. Ratios vs the hook alone are geometric means over tickets.}\label{tab:teach}
\resizebox{\columnwidth}{!}{\begin{tabular}{lrrrr}\toprule
Arm & Graph queries/run & Passed & Cost/pass & Cost vs hook\\\midrule
none & 0.28$^\dagger$ & 13 & \$0.610 & \\
hook & 0.17 (3 of 18 runs) & 14 & \$0.607 & --\\
hook + D1 teach & \textbf{1.50} (9 of 18) & 15 & \textbf{\$0.491} & $-11\%$\\
hook + D2 outline & 0.47 (4 of 17) & 12 & \$0.710 & $+1\%$\\\bottomrule
\multicolumn{5}{l}{\footnotesize $^\dagger$ the repository's own tool, which the tickets edit; no graph is built in this arm.}
\end{tabular}}\end{table}
A single line at the end of each answer (\emph{``Next time ask devkb directly: \texttt{devkb symbol X -{}-source}''})
increased direct graph queries ninefold, in half the runs, and lowered context ($-15\%$), turns ($-10\%$) and cost
(Table~\ref{tab:teach}). Agents ran the command they were shown, before editing anything. An outline added on a file's
first open did not change behaviour and added grep calls ($+18\%$). In 20 hand-checked saved answers, 20 answered the
question their intent named; one may be partial.

\subsection{Where the cost goes}\label{sec:cost}
Over 252 runs, cost correlates with turns ($r=0.92$) and context ($r=0.98$); context is 98.9\% of all tokens. A model
call costs a median \$0.017 and re-reads 41k tokens. Of those, 26.7k are present from the first call (system prompt,
tool definitions, repository instructions, ticket) and identical with or without the graph; re-read on every turn, this
fixed load is 55--68\% of a run's context. A 1{,}500-character answer is about 1\% of one turn's context. The hook's own
latency, measured on 150 recorded searches, is 0.1\,s for a pass-through and 0.8\,s for an answer, against about 17\,s
per model call. Three consequences: an answer is valuable to the extent it removes a turn; shrinking the fixed load
saves more than shrinking answers (each 1k tokens removed is about 2.5\% of every turn); and interception adds no fixed
load, unlike tool definitions.

\section{Measuring agents is itself error-prone}\label{sec:measure}
We report the errors we made, because each would have changed a conclusion. \textbf{Over-specified acceptance checks.}
Six of 24 hidden checks asked for more than their tickets stated (an implementation detail, a marker string, one of two
outcomes the ticket allowed); runs that solved the ticket scored 0 of 9. Each check was re-read against the ticket,
corrected, shown to fail on the base commit and pass on a reference fix, and every run re-judged. \textbf{Counting.} The
client's final message reports the last turn only: turns read 1 and output tokens 59 where the session made 15 calls
and 4{,}849 tokens. Per-message usage in the stream is partial too. A first rule that counted every invocation of the
tool's own binary as a self-test hid the D1 effect entirely (0.00 queries in every arm). \textbf{Concurrency and load.}
Two benchmark processes building the same checkout template corrupted 12 runs; running two batches at once doubled
suite times and made 4 of 9 suites fail that pass alone. Each is fixed in the harness (cross-process locks, per-run error
records, a re-run of failed tests with flaky tests reported apart).

\section{Threats to validity}\label{sec:threats}
\textbf{Exploratory.} The hypotheses were formed during the work; nothing here was pre-registered. \textbf{One model
and agent} (Claude Code, Sonnet). \textbf{Self-hosting}: the tickets are the tool's own, so agents may use the tool to
test their change; we separate navigation from self-tests by the first edit, a rule we had to correct once.
\textbf{Small samples}: three runs per arm; D1 rests on 6 tickets of \$0.25--0.85 each. \textbf{Moving treatment}: the
main arm's hook changed between batches as the release progressed. \textbf{Author}: one author built the tool, the
tickets, the checks and the benchmark. \textbf{Quality} is a hidden check plus a suite, not review effort,
maintainability or defects escaped.

\section{Planned confirmatory study}\label{sec:plan}
We will pre-register: \textbf{H1} intent-preserving interception answers the search's question in at least 95\% of
replayed searches; \textbf{H2} it reduces turns per successful task against native tools; \textbf{H3} the teach line
increases direct graph queries; \textbf{H4} the full configuration lowers cost per successful task; \textbf{H5} without
lowering the pass rate (H4 and H5 are tested jointly). Design: at least 100 held-out tasks from 8--12 repositories in
Python, TypeScript, Go and Java (the open-source registry of~\cite{le2026precomputed} as a start), five runs per arm,
paired by task, a mixed-effects model with task as a random effect, effect sizes with 95\% intervals; arms: native
tools; graph tools only; interception; interception + teach; at least two models and two agents. Acceptance checks are
reviewed against the task text before any run.

\section{Conclusion}
Delivering code intelligence inside the search an agent already makes improved task success without adding cost, once
the hook answered the question each search asked. Telling the agent, inside that answer, which command gave it changed
behaviour more than any instruction we had placed up front, and cut cost per successful run by a fifth in a small
sample. Because a fixed load dominates every turn, the next gains are in removing turns and shrinking that load, not in
smaller answers. These are exploratory findings; the study above is designed to confirm or refute them.

\section*{Artifact}
\ifanon Withheld for review.\else The tool, the hidden checks, the harness, per-run results and the analysis scripts are at
\url{https://github.com/namiq-ai/devkb} (\texttt{bench/results/2026-11-19-release-0.2.13/}, \texttt{docs/analysis/release-0.2.13.md}).\fi

\bibliographystyle{IEEEtran}
\bibliography{refs}
\end{document}
